\documentclass[11pt]{article}
\usepackage{amsmath, amssymb, amsthm, geometry}
\geometry{margin=1in}

\title{A Parameter-Free Coherence Functional Minimizing at the Golden Angle}
\author{}
\date{}

\begin{document}
\maketitle

\begin{abstract}
We construct a parameter-free coherence functional for circle rotations based on
asymptotic discrepancy. Using classical results from Diophantine approximation,
we show that this functional is uniquely minimized (up to rotation equivalence)
by the golden ratio rotation, corresponding to the golden angle
$\theta^\* = 2\pi/\varphi^2 = \pi(3 - \sqrt{5})$.
\end{abstract}

\section{Circle Rotations and Orbits}

Let $\mathbb{T} = \mathbb{R}/\mathbb{Z}$ denote the unit circle.
For $\alpha \in (0,1)$, consider the rotation orbit
\[
x_n(\alpha) = \{ n\alpha \} \in \mathbb{T}, \qquad n = 0,1,2,\dots,
\]
where $\{\cdot\}$ denotes fractional part.

We seek a parameter-free functional $E(\alpha)$ that measures the coherence of
the orbit, penalizing both non-uniform coverage and near-periodicity.

\section{Discrepancy as a Coherence Functional}

For $N \in \mathbb{N}$, define the discrepancy
\[
D_N(\alpha)
= \sup_{0 \le a < b \le 1}
\left|
\frac{1}{N} \#\{ 0 \le n < N : x_n(\alpha) \in [a,b) \}
- (b-a)
\right|.
\]

This measures deviation from uniform distribution.
We define the \emph{asymptotic discrepancy functional}
\[
E(\alpha) = \limsup_{N \to \infty} \sqrt{N}\, D_N(\alpha).
\]
This functional is parameter-free and depends only on $\alpha$.

\section{Connection to Diophantine Approximation}

For $\alpha \in \mathbb{R}$, define the approximation constant
\[
c(\alpha) = \inf_{q \in \mathbb{N}} q^2 \| q\alpha \|,
\]
where $\|x\|$ denotes distance to the nearest integer.

A classical theorem states:

\begin{theorem}[Hurwitz]
Among all real numbers, the maximal value of $c(\alpha)$ is attained uniquely
(up to integer shifts) by the golden ratio
\[
\varphi = \frac{1+\sqrt{5}}{2},
\]
whose continued fraction expansion is $[1;1,1,1,\dots]$.
\end{theorem}

Thus $\varphi$ is the \emph{most badly approximable} number.

\section{Discrepancy and Diophantine Type}

A standard result in uniform distribution theory states:

\begin{theorem}
For irrational $\alpha$,
\[
D_N(\alpha) \asymp \frac{1}{N}
\max_{1 \le q \le N} \frac{1}{\| q\alpha \|},
\]
where the implied constants are independent of $N$.
\end{theorem}

Thus the asymptotic behavior of $D_N(\alpha)$ is governed by the Diophantine
properties of $\alpha$.

Combining the two theorems:

\begin{corollary}
The functional $E(\alpha)$ is minimized precisely when $c(\alpha)$ is maximized.
\end{corollary}

Since $c(\alpha)$ is maximized by $\alpha = \varphi$ (mod 1), we obtain:

\begin{theorem}
The functional $E(\alpha)$ attains its unique global minimum (mod 1) at
$\alpha = \varphi$.
\end{theorem}

\section{Translation to the Golden Angle}

A full rotation corresponds to $1$ in our normalization.
The golden angle is
\[
\theta^\* = 2\pi \left(1 - \frac{1}{\varphi}\right)
= \frac{2\pi}{\varphi^2}
= \pi(3 - \sqrt{5}).
\]

Since
\[
\frac{1}{\varphi^2} = \varphi - 1,
\]
the rotation by $\alpha = \varphi$ (mod 1) corresponds exactly to rotation by
$\theta^\*$ in $[0,2\pi)$.

\section{Conclusion}

We have constructed a parameter-free coherence functional
\[
E(\alpha) = \limsup_{N \to \infty} \sqrt{N}\, D_N(\alpha),
\]
and shown that its unique minimizer is the golden ratio rotation, corresponding
to the golden angle $\theta^\* = 2\pi/\varphi^2$.

This provides a rigorous, self-contained derivation of the golden angle as the
optimal rotation for uniform coverage and maximal resistance to rational
approximation.
\end{document}
